\documentclass[../../main.tex]{subfiles}
\graphicspath{{../../image/}{../../figure/}}
\begin{document}

\section{横波相互作用}\label{section:ch15-横波相互作用}

此处的目标是估计\eqref{align:ch15-15421}右端第一项，它解释不同族黏性波之间的相互作用。在假设\eqref{align:ch15-1523}对 $t\in(0,T)$ 成立（从而\eqref{align:ch15-15420}对 $t\in[\bar t,T)$ 成立）下，将证明\eqref{align:ch15-1527}蕴含
\begin{align}
\int_{\bar t}^{T}\int_{-\infty}^{\infty}\sum_{i\neq k}\bigl(|w_iw_k|+|\omega_i\omega_k|+|w_i\omega_k|+|w_i\partial_xw_k|+|w_i\partial_x\omega_k|+|\omega_i\partial_xw_k|\bigr)\,dx\,dt=O(\delta_0^2).
\label{align:ch15-1551}
\end{align}

为此，我们比较两个抛物方程的解：
\begin{align}
\begin{cases}
\partial_tu^\flat(x,t)+\partial_x\left[\sigma^\flat(x,t)u^\flat(x,t)\right]-\partial_x^2u^\flat(x,t)=p^\flat(x,t),\\
\partial_tu^\sharp(x,t)+\partial_x\left[\sigma^\sharp(x,t)u^\sharp(x,t)\right]-\partial_x^2u^\sharp(x,t)=p^\sharp(x,t),
\end{cases}
\label{align:ch15-1552}
\end{align}
其漂移严格分离：
\begin{align}
\inf\sigma^\sharp-\sup\sigma^\flat\geqslant r>0.
\label{align:ch15-1553}
\end{align}

\begin{lemma}\label{lemma:ch15-横波相互作用估计}
若 $(u^\flat,u^\sharp)$ 是\eqref{align:ch15-1552}在 $(-\infty,\infty)\times[0,T)$ 上的解，则
\begin{align}
\int_0^T\int_{-\infty}^{\infty}|u^\flat(x,t)|\,|u^\sharp(x,t)|\,dx\,dt&\leqslant\frac1r\left\{\int_{-\infty}^{\infty}|u^\flat(x,0)|\,dx+\int_0^T\int_{-\infty}^{\infty}|p^\flat(x,t)|\,dx\,dt\right\}\nonumber\\
&\quad\times\left\{\int_{-\infty}^{\infty}|u^\sharp(x,0)|\,dx+\int_0^T\int_{-\infty}^{\infty}|p^\sharp(x,t)|\,dx\,dt\right\}.
\label{align:ch15-1554}
\end{align}
\end{lemma}

\begin{proof}
先考虑齐次情形 $p^\flat=p^\sharp=0$。引入相互作用势
\begin{align}
q(v^\flat,v^\sharp)=\int_{-\infty}^{\infty}\int_{-\infty}^{\infty}k(x-y)|v^\flat(x)|\,|v^\sharp(y)|\,dx\,dy,
\label{align:ch15-1555}
\end{align}
其中 $v^\flat$、$v^\sharp$ 为 $L^1(-\infty,\infty)$ 中任意一对函数，
\begin{align}
k(z)=
\begin{cases}
r^{-1}&z\geqslant0,\\
r^{-1}\exp\left(\frac12rz\right)&z<0.
\end{cases}
\label{align:ch15-1556}
\end{align}
注意 $rk'-2k''$ 是原点处的 Dirac 测度。我们有
\begin{align}
\frac{\mathrm{d}}{\mathrm{d}t}q(u^\flat(\cdot,t),u^\sharp(\cdot,t))&=\frac{\mathrm{d}}{\mathrm{d}t}\int_{-\infty}^{\infty}\int_{-\infty}^{\infty}k(x-y)|u^\flat(x,t)|\,|u^\sharp(y,t)|\,dx\,dy\nonumber\\
&=\int_{-\infty}^{\infty}\int_{-\infty}^{\infty}k(x-y)\left\{\left[\partial_x^2u^\flat-\partial_x(\sigma^\flat u^\flat)\right]\operatorname{sgn}u^\flat\right\}(x,t)\,|u^\sharp(y,t)|\,dx\,dy\nonumber\\
&\quad+\int_{-\infty}^{\infty}\int_{-\infty}^{\infty}k(x-y)\left\{\left[\partial_x^2u^\sharp-\partial_x(\sigma^\sharp u^\sharp)\right]\operatorname{sgn}u^\sharp\right\}(y,t)\,|u^\flat(x,t)|\,dx\,dy\nonumber\\
&=\int_{-\infty}^{\infty}\int_{-\infty}^{\infty}k'(x-y)\left[\sigma^\flat(x,t)-\sigma^\sharp(y,t)\right]|u^\flat(x,t)|\,|u^\sharp(y,t)|\,dx\,dy\nonumber\\
&\quad+\int_{-\infty}^{\infty}\int_{-\infty}^{\infty}2k''(x-y)|u^\flat(x,t)|\,|u^\sharp(y,t)|\,dx\,dy\nonumber\\
&\leqslant-\int_{-\infty}^{\infty}\int_{-\infty}^{\infty}(rk'-2k'')(x-y)|u^\flat(x,t)|\,|u^\sharp(y,t)|\,dx\,dy\nonumber\\
&\leqslant-\int_{-\infty}^{\infty}|u^\flat(x,t)|\,|u^\sharp(x,t)|\,dx.
\label{align:ch15-1557}
\end{align}
在 $(0,T)$ 上积分\eqref{align:ch15-1557}，并回顾\eqref{align:ch15-1555}与\eqref{align:ch15-1556}，得
\begin{align}
\int_0^T\int_{-\infty}^{\infty}|u^\flat(x,t)|\,|u^\sharp(x,t)|\,dx\,dt\leqslant\frac1r\int_{-\infty}^{\infty}|u^\flat(x,0)|\,dx\int_{-\infty}^{\infty}|u^\sharp(x,0)|\,dx,
\label{align:ch15-1558}
\end{align}
即特殊情形 $p^\flat=p^\sharp=0$ 下的\eqref{align:ch15-1554}。特别地，若以 $\Gamma^\flat(x,t;y,\tau)$、$\Gamma^\sharp(x,t;y,\tau)$ 记\eqref{align:ch15-1552}（齐次形式）的 Green 函数，则对任意一对初始点 $(y,\tau)$、$(y',\tau')$，
\begin{align}
\int_{\max\{\tau,\tau'\}}^T\int_{-\infty}^{\infty}\Gamma^\flat(x,t;y,\tau)\Gamma^\sharp(x,t;y',\tau')\,dx\,dt\leqslant\frac1r.
\label{align:ch15-1559}
\end{align}

\eqref{align:ch15-1552}的解现可写为
\begin{align}
\begin{cases}
u^\flat(x,t)=\displaystyle\int_{-\infty}^{\infty}\Gamma^\flat(x,t;y,0)u^\flat(y,0)\,dy+\int_0^t\int_{-\infty}^{\infty}\Gamma^\flat(x,t;y,\tau)p^\flat(y,\tau)\,dy\,d\tau,\\[2mm]
u^\sharp(x,t)=\displaystyle\int_{-\infty}^{\infty}\Gamma^\sharp(x,t;y,0)u^\sharp(y,0)\,dy+\int_0^t\int_{-\infty}^{\infty}\Gamma^\sharp(x,t;y,\tau)p^\sharp(y,\tau)\,dy\,d\tau.
\end{cases}
\label{align:ch15-15510}
\end{align}
联合\eqref{align:ch15-1559}与\eqref{align:ch15-15510}，即得\eqref{align:ch15-1554}。证明完成。
\end{proof}

\begin{lemma}\label{lemma:ch15-横波导数估计}
假设
\begin{align}
\int_0^T\int_{-\infty}^{\infty}|p^\flat(x,t)|\,dx\,dt\leqslant\delta_0,\qquad\int_0^T\int_{-\infty}^{\infty}|p^\sharp(x,t)|\,dx\,dt\leqslant\delta_0,
\label{align:ch15-15511}\\
\|\sigma^\flat(\cdot,t)\|_{L^\infty(-\infty,\infty)}\leqslant c\delta_0,\qquad\|\partial_x\sigma^\flat(\cdot,t)\|_{L^\infty(-\infty,\infty)}\leqslant c\delta_0.
\label{align:ch15-15512}
\end{align}
设 $u^\flat$、$u^\sharp$ 是\eqref{align:ch15-1552}的解，满足
\begin{align}
\|u^\flat(\cdot,t)\|_{L^1(-\infty,\infty)}\leqslant\delta_0,\qquad\|u^\sharp(\cdot,t)\|_{L^1(-\infty,\infty)}\leqslant\delta_0,
\label{align:ch15-15513}\\
\|\partial_xu^\flat(\cdot,t)\|_{L^1(-\infty,\infty)}\leqslant c\delta_0^2,\qquad\|u^\sharp(\cdot,t)\|_{L^\infty(-\infty,\infty)}\leqslant c\delta_0^2,
\label{align:ch15-15514}
\end{align}
对一切 $t\in[0,T)$ 成立。则
\begin{align}
\int_0^T\int_{-\infty}^{\infty}|\partial_xu^\flat(x,t)|\,|u^\sharp(x,t)|\,dx\,dt=O(\delta_0^2).
\label{align:ch15-15515}
\end{align}
\end{lemma}

\begin{proof}
\eqref{align:ch15-15515}左端由下式界定：
\begin{align}
\mathscr{J}(T)=\sup\int_0^{T-\tau}\int_{-\infty}^{\infty}|\partial_xu^\flat(x,t)|\,u^\sharp(x+y,t+\tau)\,dx\,dt,
\label{align:ch15-15516}
\end{align}
其中上确界取遍一切 $(y,\tau)\in(-\infty,\infty)\times[0,T)$。

由\eqref{align:ch15-15514}，
\begin{align}
\sup\int_0^1\int_{-\infty}^{\infty}|\partial_xu^\flat(x,t)|\,u^\sharp(x+y,t+\tau)\,dx\,dt\leqslant c^2\delta_0^4.
\label{align:ch15-15517}
\end{align}

对 $t>1$，把 $\partial_xu^\flat$ 写为
\begin{align}
\partial_xu^\flat(x,t)&=\int_{-\infty}^{\infty}\partial_xg(z,1)u^\flat(x-z,t-1)\,dz\nonumber\\
&\quad+\int_0^1\int_{-\infty}^{\infty}\partial_xg(z,s)\left[p^\flat-\partial_x(\sigma^\flat u^\flat)\right](x-z,t-s)\,dz\,ds,
\label{align:ch15-15518}
\end{align}
其中 $g(x,t)=(4\pi t)^{-1/2}\exp[-x^2/4t]$ 是标准热核。因此，
\begin{align}
\int_1^{T-\tau}\int_{-\infty}^{\infty}|\partial_xu^\flat(x,t)|\,u^\sharp(x+y,t+\tau)\,dx\,dt&\leqslant\int_1^{T-\tau}\int_{-\infty}^{\infty}\int_{-\infty}^{\infty}|\partial_xg(z,1)u^\flat(x-z,t-1)u^\sharp(x+y,t+\tau)|\,dz\,dx\,dt\nonumber\\
&\quad+\int_1^{T-\tau}\int_{-\infty}^{\infty}\int_0^1\int_{-\infty}^{\infty}\|\partial_x\sigma^\flat\|_{L^\infty}|\partial_xg(z,s)u^\flat(x-z,t-s)u^\sharp(x+y,t+\tau)|\,dz\,ds\,dx\,dt\nonumber\\
&\quad+\int_1^{T-\tau}\int_{-\infty}^{\infty}\int_0^1\int_{-\infty}^{\infty}\|\sigma^\flat\|_{L^\infty}|\partial_xg(z,s)\partial_xu^\flat(x-z,t-s)u^\sharp(x+y,t+\tau)|\,dz\,ds\,dx\,dt\nonumber\\
&\quad+\int_1^{T-\tau}\int_{-\infty}^{\infty}\int_{t-1}^t\int_{-\infty}^{\infty}|\partial_xg(x-z,t-s)p^\flat(z,s)u^\sharp(x+y,t+\tau)|\,dz\,ds\,dx\,dt.
\label{align:ch15-15519}
\end{align}

联合\eqref{align:ch15-15516}、\eqref{align:ch15-15517}、\eqref{align:ch15-15519}、\eqref{align:ch15-1554}、\eqref{align:ch15-15511}、\eqref{align:ch15-15512}、\eqref{align:ch15-15513}与\eqref{align:ch15-15514}，得
\begin{align}
\mathscr{J}(T)\leqslant c^2\delta_0^4+\frac{4\delta_0^2}{\sqrt{\pi}r}+\frac{8c\delta_0^3}{\sqrt{\pi}r}+\frac{2c\delta_0}{\sqrt{\pi}}\mathscr{J}(T)+\frac{2c\delta_0^3}{\sqrt{\pi}}.
\label{align:ch15-15520}
\end{align}
当 $\delta_0$ 充分小时，\eqref{align:ch15-15520}给出 $\mathscr{J}(T)=O(\delta_0^2)$，从而\eqref{align:ch15-15515}。证明完成。
\end{proof}

我们现已为建立\eqref{align:ch15-1551}做好准备。回顾\eqref{align:ch15-1526}，在\eqref{align:ch15-1554}中取 $u^\flat=w_i$，$\sigma^\flat=\sigma_i$，$p^\flat=\varphi_i$，$u^\sharp=w_k$，$\sigma^\sharp=\sigma_k$，$p^\sharp=\varphi_k$，并把时间原点从 $t=0$ 平移到 $t=\bar t$。利用\eqref{align:ch15-1527}与原书 (15.4.22)，

\end{document}
